\begin{helper}
Let $G$ be a finite graph, let $\Delta>0$, and let
$0\le\gamma\le\Delta$.  Fix distinct pairwise nonadjacent vertices
$a,b,c$ and write $Q=\{a,b,c\}$.  For the six ordered triples
\[
\mathcal O=\{(a,b,c),(a,c,b),(b,a,c),(b,c,a),(c,a,b),(c,b,a)\},
\]
let $\Phi_t$ be the full-graph ordered chamber weight of
\noderef{SamplingTripleActivationSummedSixChamberBound}.  Suppose the seven
natural numbers $n_a,n_b,n_c,n_{ab},n_{ac},n_{bc},n_{abc}$ are the actual
adjacency-pattern counts outside $Q$:
\[
n_a=|\{w\in V(G)\setminus Q:w\sim a,\ w\not\sim b,\ w\not\sim c\}|,
\]
\[
n_b=|\{w\in V(G)\setminus Q:w\not\sim a,\ w\sim b,\ w\not\sim c\}|,\quad
n_c=|\{w\in V(G)\setminus Q:w\not\sim a,\ w\not\sim b,\ w\sim c\}|,
\]
\[
n_{ab}=|\{w\in V(G)\setminus Q:w\sim a,\ w\sim b,\ w\not\sim c\}|,
\]
\[
n_{ac}=|\{w\in V(G)\setminus Q:w\sim a,\ w\not\sim b,\ w\sim c\}|,\quad
n_{bc}=|\{w\in V(G)\setminus Q:w\not\sim a,\ w\sim b,\ w\sim c\}|,
\]
\[
n_{abc}=|\{w\in V(G)\setminus Q:w\sim a,\ w\sim b,\ w\sim c\}|.
\]
Suppose also that these counts satisfy
\[
\ell_{ab}=\frac{n_{ab}+n_{abc}}{\Delta},\quad
\ell_{ac}=\frac{n_{ac}+n_{abc}}{\Delta},\quad
\ell_{bc}=\frac{n_{bc}+n_{abc}}{\Delta},
\]
\[
n_a+n_{ab}+n_{ac}+n_{abc}=\Delta,\quad
n_b+n_{ab}+n_{bc}+n_{abc}=\Delta,\quad
n_c+n_{ac}+n_{bc}+n_{abc}=\Delta,
\]
and $1-\ell_{ab},1-\ell_{ac},1-\ell_{bc}>0$.  Then
\[
\sum_{t\in\mathcal O}\Phi_t
\le
\frac1{\Delta^3}\left[
1+\frac13\sum_{pq\in\{ab,ac,bc\}}
\left(\frac{2}{(2-\ell_{pq})(1-\ell_{pq})}-1\right)\right].
\]
\end{helper}
\begin{proof}
Pair the six ordered chambers according to their two largest vertices:
$(a,b,c)$ with $(b,a,c)$, $(a,c,b)$ with $(c,a,b)$, and $(b,c,a)$ with
$(c,b,a)$.  We first identify the three exponents in each pair directly
from the displayed definition of the full-graph weight $\Phi_t$.

For $t=(a,b,c)$, the largest-coordinate exponent counts vertices outside
$Q$ adjacent to $a$, hence is
$n_a+n_{ab}+n_{ac}+n_{abc}=\Delta$.  Its middle exponent counts vertices
outside $Q$ not adjacent to $a$ and adjacent to $b$, hence is
$n_b+n_{bc}$.  Using the $b$-degree identity,
\[
n_b+n_{bc}=\Delta-(n_{ab}+n_{abc}).
\]
Its lower exponent counts vertices outside $Q$ adjacent to $c$ and adjacent
to neither $a$ nor $b$, hence is $n_c$.  For $t=(b,a,c)$, the largest
exponent is $n_b+n_{ab}+n_{bc}+n_{abc}=\Delta$, the middle exponent is
$n_a+n_{ac}=\Delta-(n_{ab}+n_{abc})$ by the $a$-degree identity, and the
lower exponent is again $n_c$.  Thus the two orders $(a,b,c)$ and
$(b,a,c)$ form the $ab$ pair with largest exponent $\Delta$, middle
exponent $\Delta-(n_{ab}+n_{abc})$, and lower exponent $n_c$.

The same calculation for $t=(a,c,b)$ and $t=(c,a,b)$ gives the $ac$ pair.
The largest exponent is $\Delta$ in both orders.  The middle exponents are
$n_c+n_{bc}$ and $n_a+n_{ab}$, respectively, and the $c$- and $a$-degree
identities rewrite both as
\[
\Delta-(n_{ac}+n_{abc}).
\]
The lower exponent in both orders is $n_b$.  Therefore this pair has
largest exponent $\Delta$, middle exponent $\Delta-(n_{ac}+n_{abc})$, and
lower exponent $n_b$.

Finally, for $t=(b,c,a)$ and $t=(c,b,a)$, the largest exponent is again
$\Delta$ in both orders.  The middle exponents are $n_c+n_{ac}$ and
$n_b+n_{ab}$, respectively, and the $c$- and $b$-degree identities rewrite
both as
\[
\Delta-(n_{bc}+n_{abc}).
\]
The lower exponent in both orders is $n_a$.  Thus the $bc$ pair has largest
exponent $\Delta$, middle exponent $\Delta-(n_{bc}+n_{abc})$, and lower
exponent $n_a$.

Set
\[
x=1-\ell_{ab},\qquad y=1-\ell_{ac},\qquad z=1-\ell_{bc},\qquad
\tau=\frac{n_{abc}}{\Delta}.
\]
The middle-exponent complement identities required for the paired chamber
normalization are therefore
\[
N_{ab}+C_{ab}=\Delta,\quad
N_{ac}+C_{ac}=\Delta,\quad
N_{bc}+C_{bc}=\Delta,
\]
where
\[
C_{ab}=n_{ab}+n_{abc},\quad
C_{ac}=n_{ac}+n_{abc},\quad
C_{bc}=n_{bc}+n_{abc},
\]
and $N_{ab},N_{ac},N_{bc}$ are the three middle exponents identified above.
The lower-exponent identities required for the paired chamber normalization
are exactly the three additive degree identities above, rewritten as
\[
n_c+(n_{ac}+n_{abc})+(n_{bc}+n_{abc})=\Delta+n_{abc},
\]
for the $ab$ lower exponent $L_{ab}=n_c$,
\[
n_b+(n_{ab}+n_{abc})+(n_{bc}+n_{abc})=\Delta+n_{abc},
\]
for the $ac$ lower exponent $L_{ac}=n_b$, and
\[
n_a+(n_{ab}+n_{abc})+(n_{ac}+n_{abc})=\Delta+n_{abc},
\]
for the $bc$ lower exponent $L_{bc}=n_a$.  These are obtained from the
$c$-, $b$-, and $a$-degree identities by adding one extra copy of
$n_{abc}$ to both sides.  Therefore
\noderef{SamplingTriplePairedChamberParameterNormalizer} applies to this
fixed independent triple with
$C_{ab}=n_{ab}+n_{abc}$, $C_{ac}=n_{ac}+n_{abc}$,
$C_{bc}=n_{bc}+n_{abc}$, $L_{ab}=n_c$, $L_{ac}=n_b$, and $L_{bc}=n_a$.
It gives nonnegative lower parameters
$\rho_x,\rho_y,\rho_z$ and the common denominator identities
\[
1+x+\rho_x=1+y+\rho_y=1+z+\rho_z=x+y+z+\tau .
\]

Using these normalized parameters,
\noderef{SamplingTripleFullChamberIntegralEstimate} bounds the sum of the
three paired finite chamber integrals by
\[
\frac{2}{(1+x)(x+y+z+\tau)}
+\frac{2}{(1+y)(x+y+z+\tau)}
+\frac{2}{(1+z)(x+y+z+\tau)}.
\]
Then \noderef{SamplingTripleRealChamberEstimate} gives
\[
\frac{2}{(1+x)(x+y+z+\tau)}
+\frac{2}{(1+y)(x+y+z+\tau)}
+\frac{2}{(1+z)(x+y+z+\tau)}
\le
1+\frac13\sum_{r\in\{x,y,z\}}\left(\frac{2}{(1+r)r}-1\right).
\]
Substituting $x=1-\ell_{ab}$, $y=1-\ell_{ac}$, and
$z=1-\ell_{bc}$ changes $(1+r)r$ into
$(2-\ell_{pq})(1-\ell_{pq})$ for the corresponding pair $pq$.
Finally, every full ordered chamber weight $\Phi_t$ contains the common
outside factor $\Delta^{-3}$, so factoring this out from the six-order sum
gives exactly the displayed inequality.
\end{proof}
